\chapter{Conclusion}
\label{ch:conclusion}

This thesis set out to identify nine Bangla dialect classes from 5-second clips
of found media, to measure how much of the apparent accuracy comes from
recognising speakers and recording channels, and to ask whether dialect identity
is carried by the sound or by the words. The answers to the four research
questions are as follows.

\paragraph{RQ1: accuracy of pretrained speech models.} Pretrained models far
outperform hand-crafted features. On a conventional clip-level split, frozen
Whisper-small embeddings reach a macro-F1 of 0.917 and fine-tuned XLS-R-300M
0.952, against 0.727 for MFCC statistics. Frozen Whisper embeddings trained on
10\,\% of the labels outperform MFCC features trained on all of them.

\paragraph{RQ2: the effect of leakage.} A strict split built from ECAPA-TDNN
pseudo-speaker clusters, calibrated against the known sources of the Formal
clips, lowers the pretrained models to 0.793 and 0.808: 12--14 points of their
clip-level scores came from speaker and channel overlap. The ranking of
pretrained over hand-crafted features survives, but fine-tuning's advantage over
frozen embeddings does not reach significance, and channel-suppression
augmentation that perturbs pitch and tempo lowers the score to 0.748.

\paragraph{RQ3: dialect information in the words.} Off-the-shelf Bangla speech
recognition loses most of it. The best recogniser, fine-tuned on regional speech
by its publishers, reproduces 19\,\% of dialect words in human transcripts;
recognisers trained on standard Bangla reproduce 1--3\,\%, rewriting dialect
into the standard form. Matching regional word lists therefore identifies the
dialect poorly (0.360 under the strict split), and a text classifier over the
full transcript reaches 0.633---between hand-crafted and pretrained acoustics.

\paragraph{RQ4: complementarity.} The routes are complementary, modestly.
Combining the text classifier with the acoustic model raises the strict-split
score from 0.808 to 0.856 (borderline under the speaker-group bootstrap) and the
clip-level score from 0.952 to 0.967 (significant). The gain is largest on the
classes the acoustic model finds hardest.

\medskip
The overall conclusion is that, on clips this short, Bangla dialect identity lies
mostly \emph{beyond the words}: in pronunciation, rhythm and intonation that
pretrained speech models capture and that speech recognition trained on the
standard language discards. The honest headline for this corpus is 0.808 from
sound alone and 0.856 with words added, on speakers the models have never heard.
Two practical lessons extend beyond Bangla: dialect and accent scores on found
data should be reported under a speaker-disjoint protocol, because the
clip-level numbers flatter every model and some more than others; and text-based
dialect methods built on speech recognition are limited by how much dialect the
recogniser preserves, which should be measured before the classifier is built.

All code, split files, transcripts, word lists and per-clip predictions are
released, and \cref{app:repro} maps every number in this thesis to the command
that produces it.
